% GENERATED FILE. Edit canonical lessons, then run npm run generate:books.
% canonical-source-hash: fnv1a64:6e2509cb500505a3
% canonical-lessons: PA-C326-darna, PA-C326-gussa, PA-C326-hairan, PA-C326-fikar, PA-C326-sharam

\chapter{Fear and Anger}
\label{ch:pa-a2-fear-and-anger}
\hlchaptermodality{326}

\begin{chapteropening}
\textbf{By the end of this chapter:} \emph{I can say what frightened or angered someone.}

Complete the last lesson of chapter 326: I can say what frightened or angered someone.
\end{chapteropening}

\section[\texorpdfstring{\pa{ਡਰਨਾ} (ḍarnā)}{ਡਰਨਾ (ḍarnā)}]{\pa{ਡਰਨਾ} (ḍarnā) — to be afraid}

\label{lesson:PA-C326-darna}



\begin{quote}
Before the new one: say the Punjabi for to chew, then the Punjabi for an ache.
\end{quote}

\subsection*{You'll want to know: \pa{ਡਰਨਾ}}
\textbf{\pa{ਡਰਨਾ}} — \emph{ḍarnā} — \textquotedblleft{}to be afraid\textquotedblright{}.

\begin{quote}
\pa{ਬੱਚਾ} \pa{ਹਨੇਰੇ} \pa{ਤੋਂ} \pa{ਡਰਦਾ} \pa{ਹੈ।}
\end{quote}

\begin{quote}
\emph{bachchā hanere ton ḍardā hai.}
\end{quote}

\begin{quote}
The child is afraid of the dark.
\end{quote}

Past (I …): \pa{ਮੈਂ} \pa{ਰਾਤ} \pa{ਨੂੰ} \pa{ਬਹੁਤ} \pa{ਡਰਿਆ}, \emph{maiṁ rāt nū̃ bahut ḍariā}.

\begin{cousinweb}
\textbf{\pa{ਡਰਨਾ}} (\emph{ḍarnā}) is inherited from Sanskrit \emph{darati}, \textquotedblleft{}is frightened\textquotedblright{}; CDIAL lists Punjabi \emph{ḍarnā}, \textquotedblleft{}to fear\textquotedblright{}, under it (CDIAL 6190).
\end{cousinweb}

\begin{grammarlens}[title={What you've built}]
One more word for this chapter.
\end{grammarlens}

\subsection*{Your turn}
\begin{itemize}
\raggedright
  \item \emph{Say it:} \emph{ḍarnā}
  \item \emph{Say it:} \emph{ḍarnā}, once more
  \item \emph{Say it:} \emph{bachchā hanere ton ḍardā hai.} — The child is afraid of the dark.
  \item \emph{Recall:} say the Punjabi for to chew, then the Punjabi for an ache, then say \emph{ḍarnā} again
\end{itemize}

\begin{tcolorbox}[breakable,colback=teal!4,colframe=teal!35!black,title={Before you move on}]
What does \pa{ਡਰਨਾ} mean? (\textquotedblleft{}To be afraid\textquotedblright{}.) Say it once more.
\end{tcolorbox}
\section[\texorpdfstring{\pa{ਗੁੱਸਾ} (gussā)}{ਗੁੱਸਾ (gussā)}]{\pa{ਗੁੱਸਾ} (gussā) — anger}

\label{lesson:PA-C326-gussa}



\begin{quote}
Before the new one: say the Punjabi for an ache, then the Punjabi for to be afraid.
\end{quote}

\subsection*{You'll want to know: \pa{ਗੁੱਸਾ}}
\textbf{\pa{ਗੁੱਸਾ}} — \emph{gussā} — \textquotedblleft{}anger\textquotedblright{}.

\begin{quote}
\pa{ਝਗੜੇ} \pa{ਤੋਂ} \pa{ਬਾਅਦ} \pa{ਅਮਨ} \pa{ਨੂੰ} \pa{ਬਹੁਤ} \pa{ਗੁੱਸਾ} \pa{ਸੀ।}
\end{quote}

\begin{quote}
\emph{jhagṛe ton bāad Aman nū̃ bahut gussā sī.}
\end{quote}

\begin{quote}
After the quarrel Aman was very angry.
\end{quote}

A masculine noun.

\begin{cousinweb}
\textbf{\pa{ਗੁੱਸਾ}} (\emph{gussā}) is a loan, usually traced through Persian to Arabic \emph{ghuṣṣa}, \textquotedblleft{}a choking; grief\textquotedblright{}.
\end{cousinweb}

\begin{grammarlens}[title={What you've built}]
One more word for this chapter.
\end{grammarlens}

\subsection*{Your turn}
\begin{itemize}
\raggedright
  \item \emph{Say it:} \emph{gussā}
  \item \emph{Say it:} \emph{gussā}, once more
  \item \emph{Say it:} \emph{jhagṛe ton bāad Aman nū̃ bahut gussā sī.} — After the quarrel Aman was very angry.
  \item \emph{Recall:} say the Punjabi for an ache, then the Punjabi for to be afraid, then say \emph{gussā} again
  \item \emph{Write it:} copy \textbf{\pa{ਗੁੱਸਾ}} once, checking the whole word
\end{itemize}

\begin{tcolorbox}[breakable,colback=teal!4,colframe=teal!35!black,title={Before you move on}]
What does \pa{ਗੁੱਸਾ} mean? (\textquotedblleft{}Anger\textquotedblright{}.) Say it once more.
\end{tcolorbox}
\section[\texorpdfstring{\pa{ਹੈਰਾਨ} (hairān)}{ਹੈਰਾਨ (hairān)}]{\pa{ਹੈਰਾਨ} (hairān) — surprised}

\label{lesson:PA-C326-hairan}



\begin{quote}
Before the new one: say the Punjabi for to be afraid, then the Punjabi for anger.
\end{quote}

\subsection*{You'll want to know: \pa{ਹੈਰਾਨ}}
\textbf{\pa{ਹੈਰਾਨ}} — \emph{hairān} — \textquotedblleft{}surprised\textquotedblright{}.

\begin{quote}
\pa{ਮੈਂ} \pa{ਉਸ} \pa{ਖ਼ਬਰ} \pa{ਤੋਂ} \pa{ਬਹੁਤ} \pa{ਹੈਰਾਨ} \pa{ਸੀ।}
\end{quote}

\begin{quote}
\emph{maiṁ us khabar ton bahut hairān sī.}
\end{quote}

\begin{quote}
I was very surprised by that news.
\end{quote}

\begin{cousinweb}
\textbf{\pa{ਹੈਰਾਨ}} (\emph{hairān}) is a loan, usually traced through Persian to Arabic \emph{ḥayrān}, \textquotedblleft{}bewildered\textquotedblright{}.
\end{cousinweb}

\begin{grammarlens}[title={What you've built}]
One more word for this chapter.
\end{grammarlens}

\subsection*{Your turn}
\begin{itemize}
\raggedright
  \item \emph{Say it:} \emph{hairān}
  \item \emph{Say it:} \emph{hairān}, once more
  \item \emph{Say it:} \emph{maiṁ us khabar ton bahut hairān sī.} — I was very surprised by that news.
  \item \emph{Recall:} say the Punjabi for to be afraid, then the Punjabi for anger, then say \emph{hairān} again
  \item \emph{Write it:} copy \textbf{\pa{ਹੈਰਾਨ}} once, checking the whole word
\end{itemize}

\begin{tcolorbox}[breakable,colback=teal!4,colframe=teal!35!black,title={Before you move on}]
What does \pa{ਹੈਰਾਨ} mean? (\textquotedblleft{}Surprised\textquotedblright{}.) Say it once more.
\end{tcolorbox}
\section[\texorpdfstring{\pa{ਫ਼ਿਕਰ} (fikar)}{ਫ਼ਿਕਰ (fikar)}]{\pa{ਫ਼ਿਕਰ} (fikar) — worry}

\label{lesson:PA-C326-fikar}



\begin{quote}
Before the new one: say the Punjabi for anger, then the Punjabi for surprised.
\end{quote}

\subsection*{You'll want to know: \pa{ਫ਼ਿਕਰ}}
\textbf{\pa{ਫ਼ਿਕਰ}} — \emph{fikar} — \textquotedblleft{}worry\textquotedblright{}.

\begin{quote}
\pa{ਮਾਤਾ} \pa{ਜੀ} \pa{ਨੂੰ} \pa{ਮੇਰਾ} \pa{ਬਹੁਤ} \pa{ਫ਼ਿਕਰ} \pa{ਹੈ।}
\end{quote}

\begin{quote}
\emph{mātā jī nū̃ merā bahut fikar hai.}
\end{quote}

\begin{quote}
Mother worries about me a lot.
\end{quote}

A masculine noun.

\begin{cousinweb}
\textbf{\pa{ਫ਼ਿਕਰ}} (\emph{fikar}) is a loan, usually traced through Persian to Arabic \emph{fikr}, \textquotedblleft{}thought\textquotedblright{}.
\end{cousinweb}

\begin{grammarlens}[title={What you've built}]
One more word for this chapter.
\end{grammarlens}

\subsection*{Your turn}
\begin{itemize}
\raggedright
  \item \emph{Say it:} \emph{fikar}
  \item \emph{Say it:} \emph{fikar}, once more
  \item \emph{Say it:} \emph{mātā jī nū̃ merā bahut fikar hai.} — Mother worries about me a lot.
  \item \emph{Recall:} say the Punjabi for anger, then the Punjabi for surprised, then say \emph{fikar} again
  \item \emph{Write it:} copy \textbf{\pa{ਫ਼ਿਕਰ}} once, checking the whole word
\end{itemize}

\begin{tcolorbox}[breakable,colback=teal!4,colframe=teal!35!black,title={Before you move on}]
What does \pa{ਫ਼ਿਕਰ} mean? (\textquotedblleft{}Worry\textquotedblright{}.) Say it once more.
\end{tcolorbox}
\section[\texorpdfstring{\pa{ਸ਼ਰਮ} (sharam)}{ਸ਼ਰਮ (sharam)}]{\pa{ਸ਼ਰਮ} (sharam) — embarrassment}

\label{lesson:PA-C326-sharam}



\begin{quote}
Before the new one: say the Punjabi for surprised, then the Punjabi for worry.
\end{quote}

\subsection*{You'll want to know: \pa{ਸ਼ਰਮ}}
\textbf{\pa{ਸ਼ਰਮ}} — \emph{sharam} — \textquotedblleft{}embarrassment\textquotedblright{}.

\begin{quote}
\pa{ਗਾਲ੍ਹ} \pa{ਤੋਂ} \pa{ਬਾਅਦ} \pa{ਉਸ} \pa{ਨੂੰ} \pa{ਬਹੁਤ} \pa{ਸ਼ਰਮ} \pa{ਆਈ।}
\end{quote}

\begin{quote}
\emph{gālh ton bāad us nū̃ bahut sharam āī.}
\end{quote}

\begin{quote}
After the swear word he felt very ashamed.
\end{quote}

A feminine noun.

\begin{cousinweb}
\textbf{\pa{ਸ਼ਰਮ}} (\emph{sharam}) is a loan, usually traced to Persian \emph{sharm}, \textquotedblleft{}shame, modesty\textquotedblright{}.
\end{cousinweb}

\begin{grammarlens}[title={What you've built}]
One more word for this chapter.
\end{grammarlens}

\subsection*{Your turn}
\begin{itemize}
\raggedright
  \item \emph{Say it:} \emph{sharam}
  \item \emph{Say it:} \emph{sharam}, once more
  \item \emph{Say it:} \emph{gālh ton bāad us nū̃ bahut sharam āī.} — After the swear word he felt very ashamed.
  \item \emph{Recall:} say the Punjabi for surprised, then the Punjabi for worry, then say \emph{sharam} again
  \item \emph{Write it:} copy \textbf{\pa{ਸ਼ਰਮ}} once, checking the whole word
\end{itemize}

\begin{tcolorbox}[breakable,colback=teal!4,colframe=teal!35!black,title={Before you move on}]
What does \pa{ਸ਼ਰਮ} mean? (\textquotedblleft{}Embarrassment\textquotedblright{}.) Say it once more.
\end{tcolorbox}
